\par\Needspace{6\baselineskip}\medskip
\begin{center}\hypertarget{native-180-article134}{}\pdfbookmark[2]{134.}{article-134}134.\end{center}
Ex praecedentibus iam ratio percipietur, cur in methodo nostra pro $x,y$
resp. quantitates
\[
\frac{n''}{n}=P\qquad\text{atque}\qquad
2\left(\frac{n+n''}{\mathpunct{\mkern-12mu.\mkern12mu}n'}-1\right)r'^3=Q
\]
accepturi simus. Patet enim \textit{primo}, si $P$ et $Q$ tamquam cognitae spectentur,
$\delta'$ inde per aequationem
\[
a\delta'=b+\frac{c+dP}{1+P}\left(1+\frac{Q}{2r'^3}\right)
\]
determinari posse, ac dein $\delta$ et $\delta''$ per aequationes 4, 6 art. 114, quum habeatur
\[
\frac{n}{n'}=\frac{1}{1+P}\left(1+\frac{Q}{2r'^3}\right)\mathpunct{\raisebox{.45ex}{,}}\qquad
\frac{n''}{n'}=\frac{P}{1+P}\left(1+\frac{Q}{2r'^3}\right)\mathpunct{\cdot}
\]
\textit{Secundo} manifestum est, in hypothesi prima pro quantitatibus $P,Q$, quarum
valores exacte veri sunt
\hypertarget{native-181}{}
\[
\frac{\theta''}{\theta}\cdot\frac{\eta}{\eta''}\mathpunct{\raisebox{.45ex}{,}}\qquad
\frac{r'r'\theta\theta''}{rr''\eta\eta''\cos f\cos f'\cos f''}\mathpunct{\raisebox{.45ex}{,}}
\]
statim obvios esse valores approximatos $\dfrac{\theta''}{\theta}\mathpunct{\raisebox{.45ex}{,}}\;\theta\theta''$, ex qua hypothesi in determinationem
ipsius $\delta'$ et proin etiam ipsarum $\delta,\delta''$ redundabunt errores primi
ordinis, vel secundi in casu speciali pluries commemorato. Ceterum etiamsi
his conclusionibus, generaliter loquendo, tutissime fidendum sit, tamen in casu
quodam speciali vim suam perdere possunt, scilicet quoties quantitas $(0.1.2)$,
quae in genere est ordinis tertii, fortuito fit $=0$ vel tam parva, ut ad altiorem
ordinem referri debeat. Hoc evenit, quoties motus geocentricus in sphaera
coelesti prope locum medium punctum inflexionis sistit. Denique apparet, ut
methodus nostra in usum vocari possit, necessario requiri, ut motus heliocentricus
inter tres observationes non nimis magnus sit: sed haec restrictio, per
problematis complicatissimi naturam, nullo modo evitari potest, neque etiam
pro incommodo habenda est, quoniam semper in votis erit, determinationem
primam orbitae incognitae corporis coelestis novi quam primum licet suscipere.
Praeterea restrictio illa sensu satis lato accipi potest, uti exempla infra tradenda
ostendent.
